\documentclass[11pt,letterpaper]{article}

\usepackage[utf8]{inputenc}
\usepackage[spanish]{babel}
\usepackage[margin=2.5cm]{geometry}
\usepackage{amsmath,amssymb,amsthm}
\usepackage{enumitem}
\usepackage{titlesec}
\usepackage{hyperref}
\usepackage{xcolor}
\usepackage{listings}

\definecolor{unicaucaverde}{RGB}{0,90,60}
\hypersetup{colorlinks=true,linkcolor=unicaucaverde,urlcolor=unicaucaverde}

\titleformat{\section}{\normalfont\Large\bfseries\color{unicaucaverde}}{\thesection}{1em}{}
\titleformat{\subsection}{\normalfont\large\bfseries}{\thesubsection}{1em}{}

\theoremstyle{definition}
\newtheorem{definicion}{Definición}[section]
\newtheorem{ejemplo}[definicion]{Ejemplo}
\newtheorem{propiedad}[definicion]{Propiedad}

\lstdefinestyle{cstyle}{
  language=C,
  basicstyle=\ttfamily\small,
  keywordstyle=\color{unicaucaverde}\bfseries,
  commentstyle=\color{gray}\itshape,
  stringstyle=\color{orange!70!black},
  numbers=left,
  numberstyle=\tiny\color{gray},
  frame=single,
  breaklines=true,
  showstringspaces=false,
  tabsize=2,
  captionpos=b
}
\lstset{style=cstyle}

\title{\textcolor{unicaucaverde}{\Large Notas de clase --- Semana 3}\\[4pt]
  \large Funciones y modularidad}
\author{Programación --- Universidad del Cauca}
\date{2026}

\begin{document}
\maketitle
\tableofcontents

\section{Funciones: qué son y por qué modularizar}

\begin{definicion}[Función]
Una función es un bloque de código con nombre propio que realiza una tarea específica, puede
recibir datos de entrada (\emph{parámetros}) y puede devolver un resultado (\emph{valor de
retorno}). Una vez definida, se puede \emph{llamar} (invocar) tantas veces como se necesite desde
cualquier parte del programa.
\end{definicion}

\begin{definicion}[Diseño descendente / modularización]
El \emph{diseño descendente} (\emph{top-down}) consiste en dividir un problema grande en
subproblemas más pequeños y manejables, y resolver cada uno con su propia función. El programa
principal (\texttt{main}) queda como una secuencia corta de llamadas a funciones, cada una
responsable de una sola tarea bien definida.
\end{definicion}

\begin{propiedad}[Ventajas de modularizar]
\begin{itemize}[leftmargin=1.2cm]
  \item \textbf{Reutilización}: una función escrita una vez se puede llamar muchas veces, en el
        mismo programa o en otros.
  \item \textbf{Legibilidad}: un \texttt{main} de 5 llamadas a funciones bien nombradas se entiende
        más rápido que 200 líneas seguidas.
  \item \textbf{Mantenibilidad}: si hay que corregir cómo se calcula algo, se corrige en un solo
        lugar (dentro de la función), no en cada copia del cálculo.
  \item \textbf{Prueba independiente}: cada función se puede probar por separado antes de integrar
        el programa completo.
\end{itemize}
\end{propiedad}

\section{Declaración, definición y prototipo de una función}

\begin{definicion}[Prototipo de función]
El \emph{prototipo} anuncia al compilador que una función existe, cuál es su nombre, qué tipo de
dato retorna y qué parámetros recibe (con sus tipos), sin incluir el cuerpo:
\begin{lstlisting}
tipo_retorno nombre_funcion(tipo1 param1, tipo2 param2, ...);
\end{lstlisting}
Se coloca antes de \texttt{main} (o en un archivo de cabecera) cuando la función se \emph{define}
después de donde se \emph{llama}.
\end{definicion}

\begin{definicion}[Definición de función]
La \emph{definición} incluye el cuerpo completo: las instrucciones que la función ejecuta, y la
instrucción \texttt{return} que devuelve el resultado (si la función no es \texttt{void}).
\end{definicion}

\begin{ejemplo}[Función con retorno]
\begin{lstlisting}
#include <stdio.h>

double areaCirculo(double radio);   // prototipo

int main(void) {
    double r = 3.0;
    double area = areaCirculo(r);   // llamada
    printf("Area = %.2f\n", area);
    return 0;
}

double areaCirculo(double radio) {  // definicion
    return 3.14159 * radio * radio;
}
\end{lstlisting}
\texttt{main} llama a \texttt{areaCirculo(r)}: el valor de \texttt{r} se copia en el parámetro
\texttt{radio}, la función calcula y devuelve el área con \texttt{return}, y ese valor queda
disponible en \texttt{main} dentro de la variable \texttt{area}.
\end{ejemplo}

\begin{definicion}[Función \texttt{void}]
Una función declarada con tipo de retorno \texttt{void} no devuelve ningún valor: se usa cuando la
tarea de la función es \emph{hacer algo} (por ejemplo, imprimir) en vez de \emph{calcular algo}. No
lleva \texttt{return valor;} (a lo sumo un \texttt{return;} vacío para salir antes de tiempo).
\end{definicion}

\begin{ejemplo}[Función \texttt{void}]
\begin{lstlisting}
#include <stdio.h>

void imprimirTabla(int n);

int main(void) {
    imprimirTabla(5);
    return 0;
}

void imprimirTabla(int n) {
    for (int i = 1; i <= 10; i++) {
        printf("%d x %d = %d\n", n, i, n * i);
    }
}
\end{lstlisting}
\end{ejemplo}

\section{Paso de parámetros por valor}

\begin{propiedad}[Paso por valor]
En C, cuando se llama a una función, cada parámetro recibe una \emph{copia} del valor del
argumento pasado. La función trabaja sobre esa copia: cualquier cambio que haga al parámetro
dentro de su cuerpo \textbf{no se refleja} en la variable original del código que hizo la llamada.
\end{propiedad}

\begin{ejemplo}[El paso por valor no modifica el original]
\begin{lstlisting}
#include <stdio.h>

void incrementar(int x) {
    x = x + 1;              // modifica solo la copia local x
    printf("Dentro de incrementar: x = %d\n", x);
}

int main(void) {
    int numero = 5;
    incrementar(numero);
    printf("Despues de incrementar: numero = %d\n", numero);
    return 0;
}
\end{lstlisting}
Salida: \texttt{Dentro de incrementar: x = 6}, seguido de \texttt{Despues de incrementar: numero
= 5}. \texttt{numero} no cambió, porque \texttt{x} es una copia independiente que vive solo
mientras se ejecuta \texttt{incrementar}. Para que una función SÍ pueda modificar una variable del
llamador hace falta pasarle su \emph{dirección de memoria} (un puntero) --- tema de la Semana 14.
\end{ejemplo}

\section{Ámbito de las variables: locales y globales}

\begin{definicion}[Variable local]
Una variable declarada \emph{dentro} de una función (incluidos sus parámetros) es \emph{local} a
esa función: solo existe y es visible mientras la función se está ejecutando, y desaparece al
terminar. Dos funciones distintas pueden tener variables locales con el mismo nombre sin
interferir entre sí.
\end{definicion}

\begin{definicion}[Variable global]
Una variable declarada \emph{fuera} de todas las funciones (normalmente al inicio del archivo) es
\emph{global}: existe durante toda la ejecución del programa y es visible desde cualquier función
definida después de ella.
\end{definicion}

\begin{ejemplo}[Variable global compartida entre funciones]
\begin{lstlisting}
#include <stdio.h>

int contadorEventos = 0;   // variable global

void registrarEvento(void) {
    contadorEventos++;     // modifica la variable global directamente
}

int main(void) {
    registrarEvento();
    registrarEvento();
    registrarEvento();
    printf("Eventos registrados: %d\n", contadorEventos);   // imprime 3
    return 0;
}
\end{lstlisting}
A diferencia de un parámetro, una variable global no se copia: todas las funciones que la usan
leen y escriben la \emph{misma} variable.
\end{ejemplo}

\begin{propiedad}[Usar variables globales con cuidado]
Abusar de variables globales dificulta entender y depurar un programa, porque cualquier función
puede cambiarlas en cualquier momento (efectos ocultos). La práctica recomendada es preferir
variables locales y comunicar datos entre funciones explícitamente, por parámetros y valores de
retorno; reservar las variables globales para casos claramente justificados (como un contador de
eventos que de verdad debe compartirse entre muchas funciones).
\end{propiedad}

\section{Ejercicios propuestos}

\begin{enumerate}[leftmargin=1.2cm]
  \item Escribe una función \texttt{double areaCirculo(double radio)} que calcule y retorne el
        área de un círculo, y un \texttt{main} que la use para un radio leído por teclado.
  \item Escribe una función \texttt{int esPar(int n)} que retorne \texttt{1} si \texttt{n} es par y
        \texttt{0} si es impar, y úsala en un ciclo para clasificar los números del 1 al 10.
  \item Escribe una función \texttt{void imprimirTabla(int n)} (sin retorno) que imprima la tabla
        de multiplicar de \texttt{n} del 1 al 10.
  \item Escribe una función \texttt{int maximo(int a, int b, int c)} que retorne el mayor de tres
        números enteros.
  \item Escribe un programa con una variable global \texttt{contadorEventos} y una función
        \texttt{void registrarEvento(void)} que la incremente; llama a la función 5 veces desde
        \texttt{main} e imprime el total al final.
  \item Dado el código
        \texttt{void incrementar(int x) \{ x = x + 1; \}}, llamado como
        \texttt{int numero = 10; incrementar(numero);}, explica por qué \texttt{numero} sigue
        valiendo 10 después de la llamada.
\end{enumerate}

\section{Soluciones de los ejercicios propuestos}

\begin{description}

\item[\textbf{Ejercicio 1.}]
\begin{lstlisting}
#include <stdio.h>

double areaCirculo(double radio);

int main(void) {
    double r;
    printf("Ingrese el radio: ");
    scanf("%lf", &r);
    printf("Area = %.2f\n", areaCirculo(r));
    return 0;
}

double areaCirculo(double radio) {
    return 3.14159 * radio * radio;
}
\end{lstlisting}

\item[\textbf{Ejercicio 2.}]
\begin{lstlisting}
#include <stdio.h>

int esPar(int n);

int main(void) {
    for (int i = 1; i <= 10; i++) {
        if (esPar(i)) {
            printf("%d es par\n", i);
        } else {
            printf("%d es impar\n", i);
        }
    }
    return 0;
}

int esPar(int n) {
    return (n % 2 == 0) ? 1 : 0;
}
\end{lstlisting}

\item[\textbf{Ejercicio 3.}]
\begin{lstlisting}
#include <stdio.h>

void imprimirTabla(int n);

int main(void) {
    int n;
    printf("Ingrese un numero: ");
    scanf("%d", &n);
    imprimirTabla(n);
    return 0;
}

void imprimirTabla(int n) {
    for (int i = 1; i <= 10; i++) {
        printf("%d x %d = %d\n", n, i, n * i);
    }
}
\end{lstlisting}

\item[\textbf{Ejercicio 4.}]
\begin{lstlisting}
#include <stdio.h>

int maximo(int a, int b, int c);

int main(void) {
    printf("El mayor es %d\n", maximo(7, 15, 3));
    return 0;
}

int maximo(int a, int b, int c) {
    int m = a;
    if (b > m) m = b;
    if (c > m) m = c;
    return m;
}
\end{lstlisting}

\item[\textbf{Ejercicio 5.}]
\begin{lstlisting}
#include <stdio.h>

int contadorEventos = 0;

void registrarEvento(void);

int main(void) {
    for (int i = 0; i < 5; i++) {
        registrarEvento();
    }
    printf("Eventos registrados: %d\n", contadorEventos);
    return 0;
}

void registrarEvento(void) {
    contadorEventos++;
}
\end{lstlisting}
Imprime \texttt{Eventos registrados: 5}.

\item[\textbf{Ejercicio 6.}]
En C, los parámetros se pasan \textbf{por valor}: al llamar a \texttt{incrementar(numero)}, el
parámetro \texttt{x} recibe una \emph{copia} del valor de \texttt{numero} (10). Dentro de la
función, \texttt{x = x + 1} modifica esa copia local (\texttt{x} pasa a valer 11), pero
\texttt{numero}, la variable original en \texttt{main}, es una variable distinta en otra región de
memoria y no se ve afectada. Al terminar la función, la copia \texttt{x} desaparece y
\texttt{numero} sigue valiendo 10.
\[\boxed{\text{\texttt{numero} sigue valiendo 10: la función solo modificó su copia local \texttt{x}.}}\]

\end{description}

\end{document}
